\documentclass[11pt]{article}
\usepackage{docstyle}
\begin{document}
\sheethead{Lesson 3 \textperiodcentered\ Homework}{NCEA Level 2 \textperiodcentered\ AS 91262 \textperiodcentered\ Sketching gradient functions}{about 60 minutes}{}
Sketch on the empty axes. Line up every key $x$-value with a ruler, and label turning points.
[A] Achieved, [M] Merit, [E] Excellence.
\par\needspace{58mm}\Q The graph of $f(x)=3x-2$ is shown. Sketch $f'(x)$.\mk{A}
\begin{center}
\sketch[xmin=-3,xmax=3,ymin=-12,ymax=8,unit=7.5mm,yunit=1.80mm,ystep=4,gy=2,name={$f(x)$},curve={3*\x-2},dmin=-3,dmax=3]\hspace{10mm}\sketch[xmin=-3,xmax=3,ymin=-4,ymax=4,unit=7.5mm,yunit=4.50mm,ystep=2,gy=1,name={$f'(x)$},ans={3},amin=-3,amax=3]
\end{center}
\par\needspace{58mm}\Q The graph of $f(x)=4-\dfrac{1}{2}x$ is shown. Sketch $f'(x)$.\mk{A}
\begin{center}
\sketch[xmin=-4,xmax=4,ymin=-1,ymax=7,unit=7.5mm,yunit=4.50mm,ystep=2,gy=1,name={$f(x)$},curve={4-0.5*\x},dmin=-4,dmax=4]\hspace{10mm}\sketch[xmin=-4,xmax=4,ymin=-2,ymax=2,unit=7.5mm,yunit=9.00mm,ystep=1,gy=1,name={$f'(x)$},ans={-0.5},amin=-4,amax=4]
\end{center}
\par\needspace{58mm}\Q The graph of $f(x)=x^{2}-4x$ is shown. Sketch $f'(x)$.\mk{A}
\begin{center}
\sketch[xmin=-1,xmax=5,ymin=-5,ymax=6,unit=7.5mm,yunit=3.27mm,ystep=2,gy=1,name={$f(x)$},curve={\x*\x-4*\x},dmin=-1,dmax=5]\hspace{10mm}\sketch[xmin=-1,xmax=5,ymin=-6,ymax=6,unit=7.5mm,yunit=3.00mm,ystep=2,gy=1,name={$f'(x)$},ans={2*\x-4},amin=-1,amax=5]
\end{center}
\par\needspace{58mm}\Q The graph of $f(x)=3-2x-x^{2}$ is shown. Sketch $f'(x)$.\mk{A}
\begin{center}
\sketch[xmin=-4,xmax=2,ymin=-6,ymax=5,unit=7.5mm,yunit=3.27mm,ystep=2,gy=1,name={$f(x)$},curve={3-2*\x-\x*\x},dmin=-4,dmax=2]\hspace{10mm}\sketch[xmin=-4,xmax=2,ymin=-6,ymax=6,unit=7.5mm,yunit=3.00mm,ystep=2,gy=1,name={$f'(x)$},ans={-2*\x-2},amin=-4,amax=2]
\end{center}
\par\needspace{58mm}\Q The graph of $f(x)=x^{3}-12x$ is shown. Sketch $f'(x)$.\mk{A}
\begin{center}
\sketch[xmin=-4,xmax=4,ymin=-18,ymax=18,unit=7.5mm,yunit=1.00mm,ystep=6,gy=3,name={$f(x)$},curve={\x*\x*\x-12*\x},dmin=-4,dmax=4]\hspace{10mm}\sketch[xmin=-4,xmax=4,ymin=-12,ymax=30,unit=7.5mm,yunit=0.86mm,ystep=6,gy=3,name={$f'(x)$},ans={3*\x*\x-12},amin=-4,amax=4]
\end{center}
\par\needspace{58mm}\Q The graph of $f(x)=-x^{3}+3x^{2}+1$ is shown. Sketch $f'(x)$.\mk{A}
\begin{center}
\sketch[xmin=-2,xmax=4,ymin=-8,ymax=12,unit=7.5mm,yunit=1.80mm,ystep=4,gy=2,name={$f(x)$},curve={-\x*\x*\x+3*\x*\x+1},dmin=-2,dmax=4]\hspace{10mm}\sketch[xmin=-2,xmax=4,ymin=-12,ymax=6,unit=7.5mm,yunit=2.00mm,ystep=3,gy=3,name={$f'(x)$},ans={-3*\x*\x+6*\x},amin=-2,amax=4]
\end{center}
\par\needspace{58mm}\Q The graph of $f(x)=\dfrac{1}{4}x^{4}-2x^{2}$ is shown. Sketch $f'(x)$.\mk{M}
\begin{center}
\sketch[xmin=-3,xmax=3,ymin=-5,ymax=4,unit=7.5mm,yunit=4.00mm,ystep=2,gy=1,name={$f(x)$},curve={0.25*\x*\x*\x*\x-2*\x*\x},dmin=-3,dmax=3]\hspace{10mm}\sketch[xmin=-3,xmax=3,ymin=-8,ymax=8,unit=7.5mm,yunit=2.25mm,ystep=2,gy=2,name={$f'(x)$},ans={\x*\x*\x-4*\x},amin=-3,amax=3]
\end{center}
\par\needspace{58mm}\Q The graph of $y=f'(x)$ is shown. Sketch a possible graph of $f(x)$.\mk{A}
\begin{center}
\sketch[xmin=-3,xmax=3,ymin=-1,ymax=4,unit=7.5mm,yunit=7.20mm,ystep=1,gy=1,name={$f'(x)$},curve={2},dmin=-3,dmax=3]\hspace{10mm}\sketch[xmin=-3,xmax=3,ymin=-8,ymax=6,unit=7.5mm,yunit=2.57mm,ystep=2,gy=1,name={$f(x)$},ans={2*\x-1},amin=-3,amax=3]
\end{center}
\par\needspace{58mm}\Q The graph of $y=f'(x)$ is shown. Sketch a possible graph of $f(x)$.\mk{A}
\begin{center}
\sketch[xmin=-1,xmax=6,ymin=-4,ymax=3,unit=7.5mm,yunit=5.14mm,ystep=1,gy=1,name={$f'(x)$},curve={\x-3},dmin=-1,dmax=6]\hspace{10mm}\sketch[xmin=-1,xmax=6,ymin=-4,ymax=6,unit=7.5mm,yunit=3.60mm,ystep=2,gy=1,name={$f(x)$},ans={0.5*\x*\x-3*\x+2},amin=-1,amax=6]
\end{center}
\par\needspace{58mm}\Q The graph of $y=f'(x)$ is shown. Sketch a possible graph of $f(x)$.\mk{A}
\begin{center}
\sketch[xmin=-5,xmax=1,ymin=-6,ymax=6,unit=7.5mm,yunit=3.00mm,ystep=2,gy=1,name={$f'(x)$},curve={-2*\x-4},dmin=-5,dmax=1]\hspace{10mm}\sketch[xmin=-5,xmax=1,ymin=-7,ymax=4,unit=7.5mm,yunit=3.27mm,ystep=2,gy=1,name={$f(x)$},ans={-\x*\x-4*\x-1},amin=-5,amax=1]
\end{center}
\par\needspace{58mm}\Q The graph of $y=f'(x)$ is shown. Sketch a possible graph of $f(x)$.\mk{M}
\begin{center}
\sketch[xmin=-3,xmax=3,ymin=-2,ymax=8,unit=7.5mm,yunit=3.60mm,ystep=2,gy=1,name={$f'(x)$},curve={\x*\x-1},dmin=-3,dmax=3]\hspace{10mm}\sketch[xmin=-3,xmax=3,ymin=-6,ymax=6,unit=7.5mm,yunit=3.00mm,ystep=2,gy=1,name={$f(x)$},ans={\x*\x*\x/3-\x},amin=-3,amax=3]
\end{center}
\par\needspace{58mm}\Q The graph of $y=f'(x)$ is shown. Sketch a possible graph of $f(x)$ and label its maximum and minimum.\mk{M}
\begin{center}
\sketch[xmin=-5,xmax=3,ymin=-12,ymax=5,unit=7.5mm,yunit=2.12mm,ystep=4,gy=2,name={$f'(x)$},curve={-\x*\x-2*\x+3},dmin=-5,dmax=3]\hspace{10mm}\sketch[xmin=-5,xmax=3,ymin=-7,ymax=6,unit=7.5mm,yunit=2.77mm,ystep=2,gy=1,name={$f(x)$},ans={-\x*\x*\x/3-\x*\x+3*\x+3},amin=-5,amax=3]
\end{center}
\par\needspace{50mm}\Q Tama says: ``Where the graph of $f'(x)$ has its highest point, the graph of $f(x)$ has a maximum.''
Is Tama right? Explain, using $f(x)=x^{3}-3x$ (so $f'(x)=3x^{2}-3$) or an example of your own.\mk{M}
\ALine{No. A maximum of $f$ is where $f'(x)=0$ and $f'$ changes from $+$ to $-$.}
\ALine{A \tturn of $f'$ is where $f$ is steepest (a \tinflect), not a \tturn of $f$.}
\ALine{E.g.\ $f'(x)=3x^{2}-3$ has its lowest point at $x=0$, where $f$ is steepest downhill;}
\ALine{the maximum of $f$ is at $x=-1$, where $f'(-1)=0$.}

\par\needspace{70mm}\QX The graph of $f'(x)=(x-2)^{2}$ touches the $x$-axis at $x=2$ but does not cross it.
Sketch a possible graph of $f(x)$, and describe what happens to $f$ at $x=2$.\mk{E}
\begin{center}
\sketch[xmin=-1,xmax=5,ymin=-1,ymax=9,unit=7.5mm,yunit=3.60mm,ystep=2,gy=1,name={$f'(x)$},curve={(\x-2)*(\x-2)},dmin=-1,dmax=5]\hspace{10mm}\sketch[xmin=-1,xmax=5,ymin=-5,ymax=6,unit=7.5mm,yunit=3.27mm,ystep=2,gy=1,name={$f(x)$},ans={(\x-2)*(\x-2)*(\x-2)/3+1},amin=-1,amax=5]
\end{center}
\ALine{$f'(x)\ge0$ everywhere, so $f$ is always increasing.}
\ALine{At $x=2$, $f'=0$ but does not change sign: $f$ is flat there (a stationary \tinflect),}
\ALineT{not a maximum or a minimum. One possible $f(x)=\dfrac{1}{3}(x-2)^{3}+1$.}

\vfill
{\small Optional extra: StudyTime Level 2 Calculus Guide, ``Sketching Gradient Functions'' from page 19.\par}
\end{document}
